\documentclass[../main.tex]{subfiles}
\begin{document}

\ifSubfilesClassLoaded{\PageLegaleChapitre}{}

\chapter{Exposants et racines carrées}
\label{chap:exposants-et-racines}

% =====================================================
\section{Les exposants (puissances)}
% =====================================================

\subsection{Définition}

\begin{definition}
Soit $a$ un nombre réel et $n$ un entier strictement positif.
\[
a^n = \underbrace{a \times a \times \cdots \times a}_{n \text{ facteurs}}
\]
Dans $a^n$, $a$ est la \textbf{base} et $n$ est l'\textbf{exposant}.

Par convention, pour $a \neq 0$ :
\[
a^1 = a, \qquad a^0 = 1, \qquad a^{-n} = \frac{1}{a^n}.
\]
\end{definition}

\textbf{Exemples.}
\[
2^5 = 2\times2\times2\times2\times2 = 32, \qquad
10^4 = 10\,000, \qquad
5^0 = 1, \qquad
2^{-3} = \frac{1}{2^3} = \frac{1}{8}.
\]

\textbf{Pourquoi ces conventions ?} Elles sont imposées par la règle du quotient (voir plus bas) :
\[
\frac{2^3}{2^3} = \frac{8}{8} = 1 \quad\text{et}\quad \frac{2^3}{2^3} = 2^{3-3} = 2^0
\quad\Rightarrow\quad 2^0 = 1,
\]
\[
\frac{2^2}{2^5} = \frac{4}{32} = \frac{1}{8} \quad\text{et}\quad \frac{2^2}{2^5} = 2^{2-5} = 2^{-3}
\quad\Rightarrow\quad 2^{-3} = \frac{1}{2^3}.
\]

\subsection{Propriétés}

Pour $a, b$ réels non nuls et $m, n$ entiers :

\begin{propriete}[Règles de calcul]
\begin{enumerate}[label=\textbf{(P\arabic*)}]
  \item \textbf{Produit de même base :} $a^m \times a^n = a^{m+n}$
  \item \textbf{Quotient de même base :} $\dfrac{a^m}{a^n} = a^{m-n}$
  \item \textbf{Puissance d'une puissance :} $(a^m)^n = a^{m \times n}$
  \item \textbf{Puissance d'un produit :} $(a \times b)^n = a^n \times b^n$
  \item \textbf{Puissance d'un quotient :} $\left(\dfrac{a}{b}\right)^n = \dfrac{a^n}{b^n}$
\end{enumerate}
\end{propriete}

\textbf{Exemples.}
\begin{itemize}
  \item $2^3 \times 2^4 = 2^{3+4} = 2^7 = 128$
        \quad (en effet : $(2\times2\times2)\times(2\times2\times2\times2)$ comporte 7 facteurs $2$)
  \item $\dfrac{5^6}{5^2} = 5^{6-2} = 5^4 = 625$
  \item $(3^2)^3 = 3^{2\times3} = 3^6 = 729$
  \item $(2\times5)^3 = 2^3 \times 5^3 = 8 \times 125 = 1\,000$
  \item $\left(\dfrac{2}{3}\right)^3 = \dfrac{2^3}{3^3} = \dfrac{8}{27}$
  \item $x^5 \times x^{-2} = x^{5+(-2)} = x^3$
\end{itemize}

\subsection{Le signe d'une puissance}

\begin{propriete}[Signe]
\begin{itemize}
  \item Si l'exposant est \textbf{pair}, la puissance d'un nombre négatif est \textbf{positive} : $(-2)^4 = 16$.
  \item Si l'exposant est \textbf{impair}, elle est \textbf{négative} : $(-2)^3 = -8$.
\end{itemize}
\end{propriete}

\begin{attention}
Les parenthèses déterminent la base :
\[
(-3)^2 = (-3)\times(-3) = 9 \qquad\text{mais}\qquad -3^2 = -(3\times 3) = -9.
\]
\end{attention}

\subsection{Erreurs fréquentes}

\begin{attention}
\begin{itemize}
  \item $2^3 \neq 2 \times 3$ : on a $2^3 = 8$ et non $6$.
  \item Les règles (P1) et (P2) exigent \textbf{la même base}. Par exemple, $2^3 \times 3^2$ ne se simplifie pas en une seule puissance : $2^3\times 3^2 = 8 \times 9 = 72$.
  \item Il n'existe pas de règle pour une \textbf{somme} : $2^3 + 2^4 = 8 + 16 = 24$, et \textbf{pas} $2^7$.
  \item $(a+b)^2 \neq a^2 + b^2$. Contre-exemple : $(3+4)^2 = 49$ alors que $3^2 + 4^2 = 25$.
  \item L'exposant ne porte que sur ce qui est juste devant : $3 \times 2^2 = 12$, alors que $(3\times2)^2 = 36$.
\end{itemize}
\end{attention}

\subsection{Puissances de 10 et notation scientifique}

\[
10^3 = 1\,000, \qquad 10^2 = 100, \qquad 10^0 = 1, \qquad 10^{-1} = 0\dec 1, \qquad 10^{-3} = 0\dec 001.
\]

\begin{definition}
Un nombre est en \textbf{notation scientifique} s'il est écrit sous la forme
\[
a \times 10^n \qquad\text{avec } 1 \leq a < 10 \text{ et } n \in \mathbb{Z}.
\]
\end{definition}

\textbf{Exemples.}
\[
45\,000 = 4\dec 5 \times 10^4, \qquad 0\dec 0072 = 7\dec 2 \times 10^{-3}.
\]

% =====================================================
\section{Les racines carrées}
% =====================================================

\subsection{Définition}

\begin{definition}
Soit $a \geq 0$. La \textbf{racine carrée} de $a$, notée $\sqrt{a}$, est l'unique nombre \textbf{positif ou nul} dont le carré vaut $a$ :
\[
\sqrt{a} \geq 0 \quad\text{et}\quad \left(\sqrt{a}\right)^2 = a.
\]
\end{definition}

\textbf{Exemples.}
\[
\sqrt{49} = 7 \text{ car } 7^2 = 49, \qquad \sqrt{0} = 0, \qquad \sqrt{\tfrac{4}{9}} = \tfrac{2}{3}, \qquad \sqrt{0\dec 25} = 0\dec 5.
\]

\textbf{Remarques.}
\begin{itemize}
  \item $(-7)^2 = 49$ également, mais $\sqrt{49} = 7$ seulement : une racine carrée n'est jamais négative.
  \item On ne peut pas calculer la racine carrée d'un nombre négatif (dans $\mathbb{R}$).
  \item Si $a \geq 0$ : $\sqrt{a^2} = a$. Pour tout réel $x$ : $\sqrt{x^2} = |x|$.
\end{itemize}

\textbf{Carrés parfaits à connaître :}
\[
\begin{array}{llllllll}
1^2 = 1 & 2^2 = 4 & 3^2 = 9 & 4^2 = 16 & 5^2 = 25 & 6^2 = 36 & 7^2 = 49 & 8^2 = 64 \\
9^2 = 81 & 10^2 = 100 & 11^2 = 121 & 12^2 = 144 & 13^2 = 169 & 14^2 = 196 & 15^2 = 225 &
\end{array}
\]

\subsection{Propriétés}

Pour $a, b \geq 0$ :

\begin{propriete}[Règles de calcul]
\begin{enumerate}[label=\textbf{(R\arabic*)}]
  \item $\left(\sqrt{a}\right)^2 = a$
  \item $\sqrt{a \times b} = \sqrt{a} \times \sqrt{b}$
  \item $\sqrt{\dfrac{a}{b}} = \dfrac{\sqrt{a}}{\sqrt{b}}$ \quad (avec $b > 0$)
\end{enumerate}
\end{propriete}

\textbf{Exemples.}
\[
\sqrt{9 \times 16} = \sqrt{144} = 12 = 3 \times 4 = \sqrt{9}\times\sqrt{16}, \qquad
\sqrt{\frac{25}{4}} = \frac{5}{2}, \qquad
\left(\sqrt{7}\right)^2 = 7.
\]

\begin{attention}
En général, $\sqrt{a+b} \neq \sqrt{a} + \sqrt{b}$ (et de même pour la soustraction).
\[
\sqrt{9+16} = \sqrt{25} = 5 \qquad\text{mais}\qquad \sqrt{9}+\sqrt{16} = 3+4 = 7.
\]
\end{attention}

\subsection{Simplifier une racine carrée}

\begin{methode}
\begin{enumerate}
  \item Chercher le \textbf{plus grand carré parfait} qui divise le nombre sous la racine.
  \item Écrire le nombre comme produit de ce carré parfait et d'un reste, puis appliquer (R2).
  \item Calculer la racine du carré parfait.
\end{enumerate}
\end{methode}

\textbf{Exemples.}
\[
\sqrt{72} = \sqrt{36 \times 2} = \sqrt{36}\times\sqrt{2} = 6\sqrt{2}, \qquad
\sqrt{50} = \sqrt{25\times 2} = 5\sqrt{2}, \qquad
\sqrt{12} = \sqrt{4 \times 3} = 2\sqrt{3}.
\]
(Ici, $6\sqrt{2}$ signifie $6 \times \sqrt{2}$.)

\subsection{Additionner et multiplier des racines}

On ne peut additionner (ou soustraire) que des racines \textbf{identiques}, exactement comme on regroupe des termes semblables ($3x + 5x = 8x$) :
\[
3\sqrt{2} + 5\sqrt{2} = 8\sqrt{2}, \qquad 7\sqrt{2} - 3\sqrt{2} = 4\sqrt{2}.
\]
L'expression $\sqrt{2}+\sqrt{3}$ ne se simplifie pas. Parfois, il faut d'abord simplifier chaque racine :
\[
\sqrt{8} + \sqrt{18} = 2\sqrt{2} + 3\sqrt{2} = 5\sqrt{2}.
\]
Pour \textbf{multiplier}, on multiplie les coefficients entre eux et les racines entre elles :
\[
2\sqrt{3} \times 3\sqrt{3} = (2\times3)\times(\sqrt{3}\times\sqrt{3}) = 6 \times 3 = 18.
\]

\subsection{Encadrer une racine carrée}

Pour situer $\sqrt{20}$ : comme $4^2 = 16 < 20 < 25 = 5^2$, on a $4 < \sqrt{20} < 5$ (la calculatrice donne $\sqrt{20} \approx 4\dec 47$).

\subsection{Lien avec le théorème de Pythagore}

Dans un triangle rectangle de côtés de l'angle droit $a$ et $b$ et d'hypoténuse $c$ :
\[
a^2 + b^2 = c^2 \quad\Longleftrightarrow\quad c = \sqrt{a^2 + b^2}.
\]
Exemple : $a = 3$ et $b = 4$ donnent $c = \sqrt{9+16} = \sqrt{25} = 5$.

% =====================================================
\section{Récapitulatif}
% =====================================================

\begin{tcolorbox}[colback=gray!6,colframe=gray!60!black,title={\bfseries À retenir}]
\textbf{Exposants} (bases identiques, non nulles)
\[
a^m a^n = a^{m+n}, \quad \frac{a^m}{a^n} = a^{m-n}, \quad (a^m)^n = a^{mn}, \quad (ab)^n = a^n b^n, \quad a^0 = 1, \quad a^{-n} = \frac{1}{a^n}.
\]
\textbf{Racines carrées} ($a, b \geq 0$)
\[
\sqrt{a} \geq 0, \quad \left(\sqrt{a}\right)^2 = a, \quad \sqrt{ab} = \sqrt{a}\sqrt{b}, \quad \sqrt{\frac{a}{b}} = \frac{\sqrt{a}}{\sqrt{b}}, \quad \sqrt{a+b} \neq \sqrt{a}+\sqrt{b}.
\]
\textbf{Réflexes :} vérifier que les bases sont identiques~; ne jamais « simplifier » une somme avec les règles du produit~; pour simplifier une racine, chercher le plus grand carré parfait.
\end{tcolorbox}

\newpage
% =====================================================
\section{Exercices}
% =====================================================

\subsection{Série A -- Briefing au sol : les exposants}

\begin{enumerate}[label=\textbf{\arabic*.}]

\item Avant le décollage, simplifie les calculs de la check-list. Écris sous la forme d'une seule puissance :
\begin{enumerate}[label=\alph*)]
  \item $3^2 \times 3^4$ \qquad
  \item $\dfrac{5^7}{5^3}$ \qquad
  \item $(2^3)^2$ \qquad
  \item $7\times7\times7\times7\times7$ \qquad
  \item $\dfrac{10^2 \times 10^5}{10^3}$
\end{enumerate}

\item Calcule :
\begin{enumerate}[label=\alph*)]
  \item $2^5$ \qquad
  \item $(-2)^3$ \qquad
  \item $-2^4$ \qquad
  \item $(-2)^4$ \qquad
  \item $10^4$ \qquad
  \item $5^0$ \qquad
  \item $2^{-3}$ \qquad
  \item $10^{-2}$
\end{enumerate}

\item Simplifie ces expressions du carnet de bord ($x \neq 0$) :
\begin{enumerate}[label=\alph*)]
  \item $x^3 \times x^5$ \qquad
  \item $\dfrac{x^9}{x^4}$ \qquad
  \item $(x^2)^4$ \qquad
  \item $(2x)^3$ \qquad
  \item $(3x^2)^2$ \qquad
  \item $x^5 \times x^{-2}$
\end{enumerate}

\item \textbf{Données de vol.} Écris en notation scientifique :
\begin{enumerate}[label=\alph*)]
  \item l'altitude de croisière : $11\,000$ m \qquad
  \item la vitesse de croisière : $900$ km/h \qquad
  \item la distance Bruxelles--New York : $5\,900$ km \qquad
  \item l'épaisseur d'un hublot : $0\dec 00045$ km
\end{enumerate}
Écris ensuite sous forme décimale :
\begin{enumerate}[label=\alph*),resume*]
  \item la vitesse du son : $3\dec 4 \times 10^2$ m/s \qquad
  \item l'altitude d'un ballon-sonde : $2\dec 5 \times 10^4$ m
\end{enumerate}

\item \textbf{Vrai ou faux ?} Corrige les affirmations fausses.
\begin{enumerate}[label=\alph*)]
  \item $2^3 \times 2^2 = 2^6$ \qquad
  \item $(3^2)^3 = 3^5$ \qquad
  \item $2^3 + 2^3 = 2^6$ \qquad
  \item $(a+b)^2 = a^2 + b^2$ \qquad
  \item $4^0 = 0$ \qquad
  \item $(-5)^2 = -25$ \qquad
  \item $10^3 \times 10^{-3} = 1$
\end{enumerate}

\end{enumerate}

\subsection{Série B -- Carnet de campagne : les racines carrées}

\begin{enumerate}[label=\textbf{\arabic*.},start=6]

\item Calcule sans calculatrice :
\begin{enumerate}[label=\alph*)]
  \item $\sqrt{16}$ \qquad
  \item $\sqrt{81}$ \qquad
  \item $\sqrt{121}$ \qquad
  \item $\sqrt{400}$ \qquad
  \item $\sqrt{0\dec 25}$ \qquad
  \item $\sqrt{\dfrac{4}{9}}$
\end{enumerate}

\item Un éclaireur repère son chemin : entre quels deux entiers consécutifs se trouve chaque racine ?
\begin{enumerate}[label=\alph*)]
  \item $\sqrt{7}$ \qquad
  \item $\sqrt{30}$ \qquad
  \item $\sqrt{80}$ \qquad
  \item $\sqrt{150}$
\end{enumerate}

\item Simplifie :
\begin{enumerate}[label=\alph*)]
  \item $\sqrt{12}$ \qquad
  \item $\sqrt{50}$ \qquad
  \item $\sqrt{48}$ \qquad
  \item $\sqrt{98}$ \qquad
  \item $\sqrt{180}$
\end{enumerate}

\item Calcule et simplifie :
\begin{enumerate}[label=\alph*)]
  \item $\sqrt{3}\times\sqrt{12}$ \qquad
  \item $4\sqrt{5} + 2\sqrt{5}$ \qquad
  \item $7\sqrt{2} - 3\sqrt{2}$ \qquad
  \item $\left(\sqrt{7}\right)^2$
  \item $3\sqrt{2} \times 2\sqrt{2}$ \qquad
  \item $\sqrt{50} - \sqrt{8}$ \qquad
  \item $\sqrt{27} + \sqrt{12}$
\end{enumerate}

\item \textbf{Vrai ou faux ?} Corrige les affirmations fausses.
\begin{enumerate}[label=\alph*)]
  \item $\sqrt{9+16} = \sqrt{9}+\sqrt{16}$ \qquad
  \item $\sqrt{9\times16} = \sqrt{9}\times\sqrt{16}$ \qquad
  \item $\sqrt{25} = \pm 5$ \qquad
  \item $2\sqrt{3} = \sqrt{6}$
\end{enumerate}

\end{enumerate}

\subsection{Série C -- Missions de vol et de camp}

\textbf{Mission 1 -- L'horizon du pilote.}
À l'altitude $h$ (en mètres), la distance à l'horizon est approximativement
\[
d \approx 3\dec 57 \times \sqrt{h} \qquad (d \text{ en km}).
\]
\begin{enumerate}[label=\alph*)]
  \item Calcule $d$ pour $h = 100$ m, puis $h = 2\,500$ m, puis $h = 10\,000$ m.
  \item Quand l'altitude est multipliée par $100$, par combien la distance à l'horizon est-elle multipliée ? Explique pourquoi.
\end{enumerate}

\textbf{Mission 2 -- Plan de vol.}
Un avion vole $120$ km vers le nord, puis $160$ km vers l'est. Quelle est la distance en ligne droite entre son point de départ et son point d'arrivée ?

\textbf{Mission 3 -- Hike en patrouille.}
Un éclaireur marche $5$ km vers l'ouest, puis $5$ km vers le sud. À quelle distance se trouve-t-il de son point de départ ? Donne la valeur exacte simplifiée, puis une valeur approchée au centième.

\textbf{Mission 4 -- Chaîne d'alerte.}
À chaque ronde, chaque personne prévenue en prévient $2$ nouvelles : à la ronde $n$, il y a donc $2^n$ nouvelles personnes prévenues.
\begin{enumerate}[label=\alph*)]
  \item Combien de personnes sont prévenues lors de la 5\textsuperscript{e} ronde ?
  \item À partir de quelle ronde dépasse-t-on $1\,000$ personnes prévenues en une seule ronde ?
\end{enumerate}

\textbf{Mission 5 -- Terrain de camp.}
La troupe dispose d'un terrain carré de $625\ \text{m}^2$. Quelle est la longueur de son côté ? Quelle longueur de corde faut-il pour en faire le tour ?

\subsection{Défi bonus}

\begin{enumerate}[label=\alph*)]
  \item $\dfrac{2^3 \times 2^5}{(2^2)^3}$ \qquad
  \item $\sqrt{2^6}$ \qquad
  \item $\sqrt{3^2 \times 5^2}$ \qquad
  \item $\sqrt{10^4}$
\end{enumerate}

\newpage
% =====================================================
\section{Corrigés}
% =====================================================

\subsection*{Série A}

\textbf{1.} a) $3^6$ \quad b) $5^4$ \quad c) $2^6$ \quad d) $7^5$ \quad e) $10^4$ (car $10^7 / 10^3$)

\textbf{2.} a) $32$ \quad b) $-8$ \quad c) $-16$ \quad d) $16$ \quad e) $10\,000$ \quad f) $1$ \quad g) $\frac{1}{8}$ \quad h) $\frac{1}{100} = 0\dec 01$

\textbf{3.} a) $x^8$ \quad b) $x^5$ \quad c) $x^8$ \quad d) $8x^3$ \quad e) $9x^4$ \quad f) $x^3$

\textbf{4.} a) $1\dec 1 \times 10^4$ m \quad b) $9 \times 10^2$ km/h \quad c) $5\dec 9 \times 10^3$ km \quad d) $4\dec 5 \times 10^{-4}$ km \quad e) $340$ m/s \quad f) $25\,000$ m

\textbf{5.}
\begin{enumerate}[label=\alph*)]
  \item Faux : $2^3 \times 2^2 = 2^5$.
  \item Faux : $(3^2)^3 = 3^6$.
  \item Faux : $2^3 + 2^3 = 8 + 8 = 16 = 2^4$ (les exposants ne s'additionnent pas dans une somme).
  \item Faux : $(a+b)^2 = a^2 + 2ab + b^2$. Contre-exemple : $(3+4)^2 = 49 \neq 25$.
  \item Faux : $4^0 = 1$.
  \item Faux : $(-5)^2 = 25$.
  \item Vrai : $10^3 \times 10^{-3} = 10^0 = 1$.
\end{enumerate}

\subsection*{Série B}

\textbf{6.} a) $4$ \quad b) $9$ \quad c) $11$ \quad d) $20$ \quad e) $0\dec 5$ \quad f) $\frac{2}{3}$

\textbf{7.} a) $2 < \sqrt{7} < 3$ \quad b) $5 < \sqrt{30} < 6$ \quad c) $8 < \sqrt{80} < 9$ \quad d) $12 < \sqrt{150} < 13$

\textbf{8.} a) $2\sqrt{3}$ \quad b) $5\sqrt{2}$ \quad c) $4\sqrt{3}$ \quad d) $7\sqrt{2}$ \quad e) $6\sqrt{5}$

\textbf{9.}
\begin{enumerate}[label=\alph*)]
  \item $\sqrt{36} = 6$
  \item $6\sqrt{5}$
  \item $4\sqrt{2}$
  \item $7$
  \item $3 \times 2 \times (\sqrt{2}\times\sqrt{2}) = 6 \times 2 = 12$
  \item $5\sqrt{2} - 2\sqrt{2} = 3\sqrt{2}$
  \item $3\sqrt{3} + 2\sqrt{3} = 5\sqrt{3}$
\end{enumerate}

\textbf{10.}
\begin{enumerate}[label=\alph*)]
  \item Faux : $\sqrt{25} = 5$ alors que $\sqrt{9}+\sqrt{16} = 7$.
  \item Vrai : les deux membres valent $12$.
  \item Faux : $\sqrt{25} = 5$ (une racine carrée est positive).
  \item Faux : $2\sqrt{3} = \sqrt{4}\times\sqrt{3} = \sqrt{12}$.
\end{enumerate}

\subsection*{Série C}

\textbf{Mission 1.}
\begin{enumerate}[label=\alph*)]
  \item $h = 100$ m : $d \approx 3\dec 57 \times 10 = 35\dec 7$ km. \\
        $h = 2\,500$ m : $\sqrt{2\,500} = 50$, donc $d \approx 178\dec 5$ km. \\
        $h = 10\,000$ m : $\sqrt{10\,000} = 100$, donc $d \approx 357$ km.
  \item La distance n'est multipliée que par $10$ (de $35\dec 7$ à $357$ km), car $\sqrt{100 \times h} = \sqrt{100}\times\sqrt{h} = 10\sqrt{h}$.
\end{enumerate}

\textbf{Mission 2.} $d = \sqrt{120^2 + 160^2} = \sqrt{14\,400 + 25\,600} = \sqrt{40\,000} = 200$ km.

\textbf{Mission 3.} $d = \sqrt{5^2 + 5^2} = \sqrt{50} = 5\sqrt{2}$ km $\approx 7\dec 07$ km.

\textbf{Mission 4.} a) $2^5 = 32$ personnes. \quad b) $2^9 = 512 < 1\,000 < 1\,024 = 2^{10}$, donc dès la 10\textsuperscript{e} ronde.

\textbf{Mission 5.} Côté : $\sqrt{625} = 25$ m. Tour complet : $4 \times 25 = 100$ m de corde.

\subsection*{Défi bonus}

a) $\dfrac{2^8}{2^6} = 2^2 = 4$ \quad b) $\sqrt{2^6} = 2^3 = 8$ (car $(2^3)^2 = 2^6$) \quad c) $3 \times 5 = 15$ \quad d) $\sqrt{10^4} = 10^2 = 100$

\end{document}
